% ============================================================================
%  XYZ File Format Family --- VSEPR-SIM Reference Document
%  VSEPR-SIM Theoretical Division
%  Classification: TOP SECRET (internal development document)
%
%  Subject: Canonical specification of the .xyz, .xyza, .xyzc, .xyzf, and
%           .xyzFull file formats used by the VSEPR-SIM simulation engine;
%           including grammar, unit conventions, PBC encoding and lattice
%           mathematics, bond auto-detection, neighbor-list construction,
%           trajectory analysis (RDF, MSD, coordination), provenance hashing,
%           per-format validation, export profile connection, interoperability,
%           and full implementation mapping.
%
%  Compile: pdflatex (twice for ToC)
% ============================================================================
\documentclass[11pt, a4paper, oneside]{article}

\usepackage[top=1.0in, bottom=1.0in, left=1.15in, right=1.15in,
            headheight=14pt]{geometry}
\usepackage[utf8]{inputenc}
\usepackage[T1]{fontenc}
\usepackage{lmodern}
\usepackage{microtype}
\usepackage{amsmath, amssymb, amsthm, mathtools}
\usepackage{bm}
\usepackage{booktabs}
\usepackage{array}
\usepackage{tabularx}
\usepackage{listings}
\usepackage{xcolor}
\usepackage{enumitem}
\usepackage{fancyhdr}
\usepackage{titlesec}
\usepackage[colorlinks=true, linkcolor=black, citecolor=black,
            urlcolor=blue]{hyperref}

% --- Code style --------------------------------------------------------------
\definecolor{codegray}{rgb}{0.5,0.5,0.5}
\definecolor{codeblue}{rgb}{0.13,0.13,0.9}
\definecolor{codegreen}{rgb}{0.0,0.5,0.0}
\definecolor{codebg}{rgb}{0.97,0.97,0.97}
\lstset{
  backgroundcolor=\color{codebg},
  basicstyle=\ttfamily\small,
  keywordstyle=\color{codeblue}\bfseries,
  commentstyle=\color{codegreen}\itshape,
  numberstyle=\tiny\color{codegray},
  numbers=left, numbersep=6pt,
  breaklines=true, frame=single,
  rulecolor=\color{codegray},
  captionpos=b, showstringspaces=false,
  tabsize=2
}

% --- Headers -----------------------------------------------------------------
\pagestyle{fancy}
\fancyhf{}
\fancyhead[L]{\small\textit{XYZ File Format Family}}
\fancyhead[R]{\small\textit{VSEPR-SIM / TOP SECRET}}
\fancyfoot[C]{\thepage}
\renewcommand{\headrulewidth}{0.4pt}

% --- Section style -----------------------------------------------------------
\titleformat{\section}{\Large\bfseries}{\S\thesection}{1em}{}
\titleformat{\subsection}{\large\bfseries}{\thesubsection}{1em}{}
\titleformat{\subsubsection}{\normalsize\bfseries}{\thesubsubsection}{1em}{}

% --- Math convenience --------------------------------------------------------
\newcommand{\ang}{\text{\AA}}
\newcommand{\fs}{\text{fs}}
\newcommand{\eV}{\text{eV}}
\newcommand{\amu}{\text{amu}}
\newcommand{\rcut}{r_{\mathrm{cut}}}
\newcommand{\rcov}{r_{\mathrm{cov}}}
\newcommand{\avec}{\mathbf{a}}
\newcommand{\bvec}{\mathbf{b}}
\newcommand{\cvec}{\mathbf{c}}
\newcommand{\rvec}{\mathbf{r}}
\newcommand{\svec}{\mathbf{s}}
\newcommand{\norm}[1]{\left\|#1\right\|}
\newcommand{\kB}{k_{\mathrm{B}}}

% ============================================================================
\begin{document}

\begin{titlepage}
  \centering
  \vspace*{1.2cm}
  {\large\scshape VSEPR-SIM Theoretical Division\par}
  \vspace{0.3cm}
  {\small TOP SECRET --- Internal Development Document\par}
  \vspace{1.2cm}
  {\Huge\bfseries XYZ File Format Family\par}
  \vspace{0.4cm}
  {\LARGE Canonical Reference for \texttt{.xyz}, \texttt{.xyza},
   \texttt{.xyzc},\\[0.2em]
   \texttt{.xyzf}, and \texttt{.xyzFull}\par}
  \vspace{1.0cm}
  {\normalsize
    Grammar\,|\,
    Unit conventions\,|\,
    PBC \& lattice mathematics\,|\,\\[0.3em]
    Bond detection\,|\,
    Trajectory analysis\,|\,
    Provenance hashing\,|\,\\[0.3em]
    Validation\,|\,
    Interoperability\,|\,
    Implementation mapping\par}
  \vspace{1.4cm}
  {\normalsize VSEPR-SIM Development Reference\par}
\end{titlepage}

\tableofcontents
\newpage

% ============================================================================
\section{Overview}
\label{sec:overview}
% ============================================================================

VSEPR-SIM uses a family of text-based coordinate files derived from the
standard XYZ format. All variants are line-oriented, human-readable, and
append-compatible. The family is ordered by information density:

\begin{table}[h]
\centering
\begin{tabular}{lllp{5.2cm}}
  \toprule
  Extension & Name & Per-atom data & Primary use \\
  \midrule
  \texttt{.xyz}     & Standard XYZ    & species, position
    & Import/export baseline; external tool interoperability \\
  \texttt{.xyza}    & Extended XYZ    & + charge, velocity, force
    & Production trajectory output \\
  \texttt{.xyzc}    & Checkpoint XYZ  & + full header block
    & Simulation restart; checkpoint/restart files \\
  \texttt{.xyzf}    & Formation XYZ   & species, position (multi-frame)
    & Formation trajectories; NEB images; defect snapshots \\
  \texttt{.xyzFull} & Full-State XYZ  & + stress, KE, PE, image count
    & Rich replay; analysis pipeline carrier \\
  \bottomrule
\end{tabular}
\caption{Format hierarchy.}
\label{tab:hierarchy}
\end{table}

All formats share the same per-atom coordinate block grammar. They differ
in comment-line content, per-atom column set, and optional header blocks.
A file in any format is a valid read target for any parser one level below
it in the hierarchy; extra columns are discarded gracefully.

% ----------------------------------------------------------------------------
\subsection{Unit Conventions}
\label{sec:units}
% ----------------------------------------------------------------------------

\begin{table}[h]
\centering
\begin{tabular}{llll}
  \toprule
  Quantity & Unit & Symbol & Notes \\
  \midrule
  Position              & \AA ngstr\"{o}m      & $\ang$                 & \\
  Velocity              & \AA/femtosecond      & $\ang\,\fs^{-1}$       & \\
  Force                 & eV/\AA               & $\eV\,\ang^{-1}$       & \\
  Charge                & elementary charge    & $e$                    & \\
  Mass                  & atomic mass units    & $\amu$                 & \\
  Energy (total)        & electronvolt         & $\eV$                  & per cell \\
  Energy (per atom)     & electronvolt         & $\eV\,\text{atom}^{-1}$& in \texttt{.xyzFull} \\
  Time                  & femtosecond          & $\fs$                  & \\
  Stress                & eV/\AA$^3$           & $\eV\,\ang^{-3}$       & \texttt{.xyzFull} only \\
  Temperature           & kelvin               & K                      & comment line \\
  Pressure              & gigapascal           & GPa                    & header only \\
  \bottomrule
\end{tabular}
\caption{Unit system (all formats).}
\label{tab:units}
\end{table}

Units are implicit --- never written into the file. All readers and writers
in the VSEPR-SIM codebase assume the above system without conversion.
Callers receiving data from external tools must convert before passing
frames to VSEPR-SIM APIs.

% ----------------------------------------------------------------------------
\subsection{Coordinate System}
\label{sec:coords}
% ----------------------------------------------------------------------------

All formats use a right-handed Cartesian coordinate system with the
$z$-axis as the primary symmetry axis for layered and planar systems.
The origin is not fixed: positions are absolute within the simulation box,
which may be centred at the origin or offset depending on the lattice
generator used. For periodic systems the box is defined by three lattice
vectors $\avec$, $\bvec$, $\cvec$ encoded in the comment line
(see \S\ref{sec:pbc}).

% ----------------------------------------------------------------------------
\subsection{Reserved Species Tokens}
\label{sec:reserved_tokens}
% ----------------------------------------------------------------------------

In addition to standard element symbols, the following tokens are valid
in the species column of any XYZ variant:

\begin{center}
\begin{tabular}{ll}
  \toprule
  Token & Meaning \\
  \midrule
  \texttt{alpha}   & Alpha particle ($^4$He$^{2+}$ proxy) \\
  \texttt{bead}    & Coarse-grained polymer/bead particle \\
  \texttt{ghost}   & Non-interacting probe particle \\
  \texttt{vacancy} & Explicit vacancy marker (no mass, no force) \\
  \texttt{dummy}   & Geometry placeholder; excluded from energy \\
  \bottomrule
\end{tabular}
\end{center}

Reserved tokens are case-sensitive. They resolve to zero mass and zero
covalent radius unless overridden in \texttt{element\_data.hpp}.

% ----------------------------------------------------------------------------
\subsection{Connection to the Artifact Stack}
\label{sec:artifact_stack}
% ----------------------------------------------------------------------------

The XYZ family integrates with the wider VSEPR-SIM output stack defined
in WO-CHART-FIELD-B:

\begin{center}
\begin{tabular}{lp{8cm}}
  \toprule
  XYZ format & Downstream consumers \\
  \midrule
  \texttt{.xyza}    & Analysis pipeline: RDF, MSD, coordination number,
                       field projection; also \texttt{.xchart} trajectory views \\
  \texttt{.xyzc}    & Simulation restart; checkpoint integrity hash \\
  \texttt{.xyzf}    & Formation event timeline in \texttt{.dynx};
                       NEB chart projections in \texttt{.xchart} \\
  \texttt{.xyzFull} & Rich replay in Qt desktop; full analysis pipeline;
                       input to \texttt{.xchart} property surfaces \\
  \bottomrule
\end{tabular}
\end{center}

% ============================================================================
\section{Standard XYZ (\texttt{.xyz})}
\label{sec:xyz}
% ============================================================================

The standard XYZ format is the baseline interoperability format. It is
accepted by VESTA, Ovito, ASE, and most molecular visualisation tools
without modification.

% ----------------------------------------------------------------------------
\subsection{Grammar}
\label{sec:xyz_grammar}
% ----------------------------------------------------------------------------

\begin{lstlisting}[caption={Standard XYZ grammar (EBNF sketch).}]
xyz-file   ::= frame+
frame      ::= atom-count NEWLINE comment NEWLINE atom-line+
atom-count ::= INTEGER
comment    ::= STRING        (* free-form; parsed by VSEPR-SIM convention *)
atom-line  ::= SPECIES REAL REAL REAL NEWLINE
SPECIES    ::= element-symbol | reserved-token
\end{lstlisting}

Fields are whitespace-delimited. Tabs and multiple spaces are both accepted.
Trailing whitespace is ignored. Blank lines between frames are not permitted;
they cause a parse error. The parser is strict on frame boundaries but
lenient on internal whitespace.

% ----------------------------------------------------------------------------
\subsection{Comment Line Convention}
\label{sec:xyz_comment}
% ----------------------------------------------------------------------------

VSEPR-SIM stores structured metadata in the comment line using a
\texttt{key=value} syntax. Keys are lowercase and space-separated.
Values containing spaces must be enclosed in double quotes.

\begin{lstlisting}[caption={Full comment line example.}]
step=1500 time_fs=75.000 energy_eV=-342.819 temperature_K=301.4
  pbc="T T T" lattice="22.624 0.0 0.0  0.0 22.624 0.0  0.0 0.0 22.624"
  label=nacl_md hash=a3f82c91
\end{lstlisting}

\begin{table}[h]
\centering
\begin{tabular}{llll}
  \toprule
  Key & Type & Meaning & Formats \\
  \midrule
  \texttt{step}            & int    & Simulation step index & all \\
  \texttt{time\_fs}        & float  & Simulation time (fs)  & all \\
  \texttt{energy\_eV}      & float  & Total system energy (eV) & all \\
  \texttt{temperature\_K}  & float  & Instantaneous kinetic temperature (K) & all \\
  \texttt{pbc}             & string & \texttt{"T T T"} per-axis PBC flags & all \\
  \texttt{lattice}         & string & $3\times3$ box matrix, row-major (\AA) & all \\
  \texttt{label}           & string & Free-text run label & all \\
  \texttt{hash}            & string & Frame provenance hash (8 hex chars) & all \\
  \texttt{frame\_type}     & string & Physical frame label (see \S\ref{sec:xyzf}) & \texttt{.xyzf} \\
  \texttt{Properties}      & string & ASE column declaration & \texttt{.xyza} \\
  \texttt{ke\_eV}          & float  & Total kinetic energy (eV) & \texttt{.xyzFull} \\
  \texttt{pe\_eV}          & float  & Total potential energy (eV) & \texttt{.xyzFull} \\
  \texttt{virial\_eV}      & float  & Scalar virial (eV) & \texttt{.xyzFull} \\
  \bottomrule
\end{tabular}
\caption{Recognised comment-line keys.}
\label{tab:comment_keys}
\end{table}

Unknown keys are silently ignored, enabling forward compatibility.

% ----------------------------------------------------------------------------
\subsection{Example: Water}
\label{sec:xyz_example}
% ----------------------------------------------------------------------------

\begin{lstlisting}[caption={Minimal water molecule (\texttt{.xyz}).}]
3
step=0 time_fs=0.000 energy_eV=-14.203 label=H2O_relax
O   0.000000   0.000000   0.119748
H   0.000000   0.756950  -0.478993
H   0.000000  -0.756950  -0.478993
\end{lstlisting}

% ----------------------------------------------------------------------------
\subsection{Multi-Frame Trajectory}
\label{sec:xyz_trajectory}
% ----------------------------------------------------------------------------

Multiple frames are concatenated without separators. A reader must
consume exactly \texttt{atom\_count} atom lines per frame before beginning
the next frame header. The total number of frames is not known in advance;
readers should stream rather than pre-allocate.

Frame count given file size $F$ (bytes), average bytes per atom line
$\bar{b}$, and comment/header overhead $c$:
\begin{equation}
  N_{\mathrm{frames}} \approx \frac{F}{c + N_{\mathrm{atoms}}\,\bar{b}}
  \label{eq:frame_count}
\end{equation}

This estimate is approximate because $\bar{b}$ depends on coordinate
precision and species name length. Use it only for pre-allocation hints.

% ----------------------------------------------------------------------------
\subsection{Parsing Notes and Error Handling}
\label{sec:xyz_parsing}
% ----------------------------------------------------------------------------

\begin{itemize}
  \item The atom count on line~1 is authoritative. Readers must not count
    atom lines manually.
  \item Species tokens are case-sensitive. \texttt{Na} and \texttt{na} are
    different tokens; only \texttt{Na} is valid.
  \item Coordinates beyond the fourth column are ignored in \texttt{.xyz}
    mode. They are meaningful only in \texttt{.xyza} and \texttt{.xyzFull}.
  \item An empty comment line is valid; all comment fields default to
    zero or empty string.
  \item A mismatch between the declared atom count and the number of
    atom lines before the next frame header is a \textbf{hard error}.
    The parser throws \texttt{XyzParseError} with the offending line number.
  \item Non-finite coordinate values (\texttt{nan}, \texttt{inf}) are
    treated as a hard error. The simulation has diverged and the file
    should not be silently passed to downstream tools.
  \item Lines beginning with \texttt{\#} outside a checkpoint header block
    are treated as comments and skipped.
\end{itemize}

% ============================================================================
\section{Extended XYZ (\texttt{.xyza})}
\label{sec:xyza}
% ============================================================================

The extended format appends per-atom dynamical quantities to each atom line.
It is the primary output format for VSEPR-SIM production runs and the main
input to the analysis pipeline.

% ----------------------------------------------------------------------------
\subsection{Grammar}
\label{sec:xyza_grammar}
% ----------------------------------------------------------------------------

\begin{lstlisting}[caption={Extended XYZ atom-line grammar.}]
atom-line ::= SPECIES  X Y Z  Q  VX VY VZ  FX FY FZ  NEWLINE
              (* all numeric fields are REAL *)
\end{lstlisting}

Column order is fixed. A reader that encounters fewer than 11 data columns
per atom line shall treat all missing trailing columns as zero. A reader
encountering more than 11 columns shall ignore the excess unless it is
operating in \texttt{.xyzFull} mode.

% ----------------------------------------------------------------------------
\subsection{Property Columns}
\label{sec:xyza_cols}
% ----------------------------------------------------------------------------

\begin{table}[h]
\centering
\begin{tabular}{rllll}
  \toprule
  Col & Field & Symbol & Unit & Notes \\
  \midrule
  1   & Species    & ---  & ---                & Element or reserved token \\
  2   & $x$        & $x$  & $\ang$             & Cartesian $x$ position \\
  3   & $y$        & $y$  & $\ang$             & Cartesian $y$ position \\
  4   & $z$        & $z$  & $\ang$             & Cartesian $z$ position \\
  5   & Charge     & $q$  & $e$                & Partial atomic charge \\
  6   & $v_x$      & $v_x$& $\ang\,\fs^{-1}$   & Velocity $x$-component \\
  7   & $v_y$      & $v_y$& $\ang\,\fs^{-1}$   & Velocity $y$-component \\
  8   & $v_z$      & $v_z$& $\ang\,\fs^{-1}$   & Velocity $z$-component \\
  9   & $F_x$      & $F_x$& $\eV\,\ang^{-1}$   & Force $x$-component \\
  10  & $F_y$      & $F_y$& $\eV\,\ang^{-1}$   & Force $y$-component \\
  11  & $F_z$      & $F_z$& $\eV\,\ang^{-1}$   & Force $z$-component \\
  \bottomrule
\end{tabular}
\caption{Extended property columns (fixed order).}
\label{tab:xyza_cols}
\end{table}

% ----------------------------------------------------------------------------
\subsection{Properties Declaration Syntax}
\label{sec:xyza_props}
% ----------------------------------------------------------------------------

For interoperability with ASE extended-XYZ, VSEPR-SIM optionally writes a
\texttt{Properties} key in the comment line:

\begin{lstlisting}[caption={Properties key (ASE compatibility).}]
Properties=species:S:1:pos:R:3:charge:R:1:vel:R:3:forces:R:3
\end{lstlisting}

The value is a colon-delimited list of \texttt{name:type:count} triples.
Type codes: \texttt{S} = string, \texttt{R} = real, \texttt{I} = integer.
This key is optional; VSEPR-SIM readers use the fixed column order
regardless.

% ----------------------------------------------------------------------------
\subsection{Example: NaCl Unit Cell}
\label{sec:xyza_example}
% ----------------------------------------------------------------------------

\begin{lstlisting}[caption={Two-ion NaCl unit at 300\,K (\texttt{.xyza}).}]
2
step=200 time_fs=10.000 energy_eV=-8.142 temperature_K=300.1
  pbc="T T T" lattice="5.640 0.0 0.0  0.0 5.640 0.0  0.0 0.0 5.640"
Na  0.000000  0.000000  0.000000  +1.0  0.023 -0.011  0.008  0.312 -0.197  0.143
Cl  2.820000  2.820000  2.820000  -1.0 -0.023  0.011 -0.008 -0.312  0.197 -0.143
\end{lstlisting}

% ----------------------------------------------------------------------------
\subsection{Charge Semantics}
\label{sec:xyza_charge}
% ----------------------------------------------------------------------------

The charge column carries the value from the active charge model:

\begin{center}
\begin{tabular}{llll}
  \toprule
  Charge model & Source & Typical range & Continuity \\
  \midrule
  \texttt{formal}  & Integer valence assignment     & $\{-4,\ldots,+4\}$ & discrete \\
  \texttt{neutral} & Zero for all species           & $0$                & --- \\
  \texttt{bader}   & Bader partitioning             & $\sim\!\pm 2$      & continuous \\
  \texttt{eam}     & Not applicable (metallic)      & $0$                & --- \\
  \bottomrule
\end{tabular}
\end{center}

The active charge model is determined by the \texttt{[material]} registry
field \texttt{default\_charge\_model} (see \texttt{VSIM\_LANGUAGE.md}).
Mixing charge models across frames of the same trajectory is a validation
error (see \S\ref{sec:validation}).

% ============================================================================
\section{Checkpoint XYZ (\texttt{.xyzc})}
\label{sec:xyzc}
% ============================================================================

The checkpoint format extends \texttt{.xyza} with a structured header block
carrying all state required to resume a simulation without re-reading the
original script. A single \texttt{.xyzc} file is a complete restart package.

% ----------------------------------------------------------------------------
\subsection{Header Block}
\label{sec:xyzc_header}
% ----------------------------------------------------------------------------

The header block precedes the first frame. It begins with
\texttt{\#\# XYZC HEADER BEGIN} and ends with \texttt{\#\# XYZC HEADER END}.
All interior lines begin with \texttt{\#} and use \texttt{key = value} syntax.

\begin{lstlisting}[caption={Checkpoint header block (\texttt{.xyzc}).}]
## XYZC HEADER BEGIN
# format_version = 2
# project        = nacl_relax_v3
# vsim_version   = v5.0.0-beta.9
# seed_base      = 1001
# step           = 4000
# time_fs        = 200.000
# dt_fs          = 0.500
# run_mode       = md
# solver         = verlet
# forcefield     = ewald_formal
# temperature_K  = 300.0
# pressure_GPa   = 0.0
# periodic       = true
# pbc_flags      = T T T
# lattice_ang    = 22.624 0.000 0.000  0.000 22.624 0.000  0.000 0.000 22.624
# n_atoms        = 512
# charge_model   = formal
# label          = nacl_4x4x4_md_300K
# file_hash      = a3f82c91d48e7b02
## XYZC HEADER END
512
step=4000 time_fs=200.000 energy_eV=-17641.320 temperature_K=299.8
Na  0.000000  0.000000  0.000000  +1.0  ...
\end{lstlisting}

% ----------------------------------------------------------------------------
\subsection{Header Field Semantics}
\label{sec:xyzc_fields}
% ----------------------------------------------------------------------------

\begin{table}[h]
\centering
\begin{tabular}{llp{6cm}}
  \toprule
  Field & Type & Meaning \\
  \midrule
  \texttt{format\_version}  & int    & Header schema version (currently 2) \\
  \texttt{project}          & string & Script \texttt{[project].name} value \\
  \texttt{vsim\_version}    & string & VSEPR-SIM build version string \\
  \texttt{seed\_base}       & int    & RNG seed used for the run \\
  \texttt{step}             & int    & Last completed step index \\
  \texttt{time\_fs}         & float  & Simulation time at checkpoint (fs) \\
  \texttt{dt\_fs}           & float  & Integration timestep (fs) \\
  \texttt{run\_mode}        & string & \texttt{relax}, \texttt{md}, \texttt{npt}, \ldots \\
  \texttt{solver}           & string & Active integrator tag \\
  \texttt{forcefield}       & string & Active force-field tag \\
  \texttt{temperature\_K}   & float  & Target temperature (K) \\
  \texttt{pressure\_GPa}    & float  & Target pressure (GPa; NPT only) \\
  \texttt{periodic}         & bool   & PBC enabled at checkpoint \\
  \texttt{pbc\_flags}       & string & \texttt{"T T T"} per-axis flags \\
  \texttt{lattice\_ang}     & string & $3\times3$ lattice matrix, row-major (\AA) \\
  \texttt{n\_atoms}         & int    & Particle count (must match frame) \\
  \texttt{charge\_model}    & string & \texttt{formal}, \texttt{neutral}, \texttt{bader} \\
  \texttt{label}            & string & Free-text run identifier \\
  \texttt{file\_hash}       & string & SHA-256 prefix of the full file at write time \\
  \bottomrule
\end{tabular}
\caption{Checkpoint header fields.}
\label{tab:xyzc_header}
\end{table}

All fields are written by VSEPR-SIM. A reader encountering an absent
field falls back to the current script's values.

% ----------------------------------------------------------------------------
\subsection{Restart Protocol}
\label{sec:xyzc_restart}
% ----------------------------------------------------------------------------

\begin{enumerate}
  \item Parse the \texttt{\#\# XYZC HEADER} block into a
    \texttt{CheckpointHeader} struct.
  \item Validate \texttt{n\_atoms} against the subsequent frame's atom
    count. Mismatch is a hard error.
  \item Validate the \texttt{file\_hash} field against a freshly computed
    hash of the file's first $10\,\mathrm{KB}$ (integrity check).
  \item Override the active script's \texttt{[run]} and \texttt{[simulation]}
    fields with the header values: step counter, timestep, solver,
    force-field, PBC flags, and lattice.
  \item Load positions, velocities, forces, and charges directly into the
    integrator's state arrays.
  \item Resume the time loop from \texttt{step + 1}.
\end{enumerate}

\noindent\textbf{Note.} Checkpoint restart does not re-run the presolve or
lattice-builder phases. Geometry is taken entirely from the file. If the
original lattice-builder would produce a different geometry (e.g.\ due to
a seed change), the checkpoint takes priority.

% ----------------------------------------------------------------------------
\subsection{Checkpoint vs.\ Trajectory}
\label{sec:xyzc_vs_traj}
% ----------------------------------------------------------------------------

A \texttt{.xyzc} file written every $N$ steps doubles as a trajectory: the
header block appears once at the start; subsequent frames are appended as
plain \texttt{.xyza} blocks. To seek to frame $k$, a reader computes the
line offset:

\begin{equation}
  \mathrm{offset}(k)
  = N_{\mathrm{header}} + k \times (N_{\mathrm{atoms}} + 2)
  \label{eq:frame_offset}
\end{equation}

where $N_{\mathrm{header}}$ is the number of header lines (including the
two sentinel lines), and the $+2$ accounts for the atom-count line and the
comment line per frame. This is deterministic given \texttt{n\_atoms}.

% ============================================================================
\section{Formation XYZ (\texttt{.xyzf})}
\label{sec:xyzf}
% ============================================================================

The formation format records the positional trajectory of a formation or
reaction event as a sequence of standard coordinate frames. It omits
velocities and forces to keep file size minimal.

% ----------------------------------------------------------------------------
\subsection{Grammar}
\label{sec:xyzf_grammar}
% ----------------------------------------------------------------------------

\begin{lstlisting}[caption={Formation XYZ grammar.}]
xyzf-file  ::= frame+
frame      ::= atom-count NEWLINE comment NEWLINE atom-line+
atom-line  ::= SPECIES  X Y Z  NEWLINE
\end{lstlisting}

The grammar is syntactically identical to \texttt{.xyz}. The semantic
distinction is that the comment line always carries a \texttt{frame\_type}
key indicating the physical role of the frame.

% ----------------------------------------------------------------------------
\subsection{Comment Line Convention}
\label{sec:xyzf_comment}
% ----------------------------------------------------------------------------

\begin{lstlisting}[caption={Formation XYZ comment lines over a reaction event.}]
frame=0 frame_type=pre_formation  step=0  energy_eV=-342.100
frame=1 frame_type=transition     step=45 energy_eV=-339.817
frame=2 frame_type=post_formation step=80 energy_eV=-346.203
\end{lstlisting}

\begin{center}
\begin{tabular}{lp{8cm}}
  \toprule
  \texttt{frame\_type} token & Meaning \\
  \midrule
  \texttt{pre\_formation}  & Initial, unformed state \\
  \texttt{transition}      & Intermediate or saddle-point frame \\
  \texttt{post\_formation} & Final, formed state \\
  \texttt{defect\_snap}    & Defect snapshot during a longer trajectory \\
  \texttt{neb\_image}      & Nudged-elastic-band intermediate image \\
  \texttt{arbitrary}       & No specific physical label assigned \\
  \bottomrule
\end{tabular}
\end{center}

% ----------------------------------------------------------------------------
\subsection{Reaction Coordinate}
\label{sec:xyzf_rxcoord}
% ----------------------------------------------------------------------------

For NEB pathways and reaction trajectories, VSEPR-SIM writes a
\texttt{rxcoord} key in the comment line encoding the reaction coordinate
$\xi \in [0,1]$:

\begin{equation}
  \xi_k = \frac{\sum_{j=1}^{k} \norm{\rvec^{(j)} - \rvec^{(j-1)}}}
               {\sum_{j=1}^{M} \norm{\rvec^{(j)} - \rvec^{(j-1)}}}
  \label{eq:rxcoord}
\end{equation}

where $\rvec^{(k)}$ is the concatenated position vector of all atoms at
frame $k$, and $M+1$ is the total number of frames. This maps the
discrete image sequence onto a normalised path length.

% ----------------------------------------------------------------------------
\subsection{Frame Independence and Variable Atom Counts}
\label{sec:xyzf_independence}
% ----------------------------------------------------------------------------

Each frame is fully self-contained. The atom count may differ between frames
(e.g.\ if atoms are inserted or removed during a formation event). Readers
must re-parse the atom count at the start of each frame and must not assume
count consistency. If frame-to-frame count changes are detected, a
\texttt{variable\_count} flag is emitted in the metadata.

% ----------------------------------------------------------------------------
\subsection{Typical Uses}
\label{sec:xyzf_uses}
% ----------------------------------------------------------------------------

\begin{table}[h]
\centering
\begin{tabular}{ll}
  \toprule
  Use case & Description \\
  \midrule
  Formation event recording & Capture atomic positions during bond formation \\
  Defect migration path     & Snapshot sequence along a migration trajectory \\
  NEB pathway               & Nudged-elastic-band intermediate images \\
  Visualisation export      & Compact trajectory for external viewers \\
  Batch comparison          & Multiple independent frames in one file \\
  Reaction pathway archive  & Store full reaction path with energy profile \\
  \bottomrule
\end{tabular}
\caption{Common \texttt{.xyzf} applications.}
\label{tab:xyzf_uses}
\end{table}

% ============================================================================
\section{Full-State XYZ (\texttt{.xyzFull})}
\label{sec:xyzFull}
% ============================================================================

The \texttt{.xyzFull} format is the richest member of the family. It extends
\texttt{.xyza} with per-atom stress contributions, kinetic and potential
energy decompositions, and PBC image counters. Its primary roles are rich
trajectory replay in the Qt desktop and serving as the input carrier for the
full analysis pipeline (RDF, MSD, field projection, scale sampling).

% ----------------------------------------------------------------------------
\subsection{Grammar}
\label{sec:xyzFull_grammar}
% ----------------------------------------------------------------------------

\begin{lstlisting}[caption={\texttt{.xyzFull} atom-line grammar.}]
atom-line ::= SPECIES  X Y Z  Q  VX VY VZ  FX FY FZ
              KE  PE  SXX SYY SZZ SXY SXZ SYZ
              IX IY IZ  NEWLINE
(* 22 fields per line: 1 string + 21 reals *)
\end{lstlisting}

\begin{table}[h]
\centering
\begin{tabular}{rllll}
  \toprule
  Col(s) & Field & Symbol & Unit & Notes \\
  \midrule
  1      & Species     & ---       & ---                       & \\
  2--4   & Position    & $x,y,z$   & $\ang$                    & \\
  5      & Charge      & $q$       & $e$                       & \\
  6--8   & Velocity    & $\mathbf{v}$ & $\ang\,\fs^{-1}$      & \\
  9--11  & Force       & $\mathbf{F}$ & $\eV\,\ang^{-1}$      & \\
  12     & KE per atom & $\varepsilon_k$ & $\eV$              & $\frac{1}{2}mv^2$ \\
  13     & PE per atom & $\varepsilon_p$ & $\eV$              & force-field derived \\
  14--19 & Stress tensor & $\sigma_{\alpha\beta}$ & $\eV\,\ang^{-3}$ & Voigt order: $xx,yy,zz,xy,xz,yz$ \\
  20--22 & Image count & $n_x,n_y,n_z$ & ---                 & PBC image displacement \\
  \bottomrule
\end{tabular}
\caption{\texttt{.xyzFull} property columns (fixed order).}
\label{tab:xyzFull_cols}
\end{table}

% ----------------------------------------------------------------------------
\subsection{Stress Tensor Convention}
\label{sec:xyzFull_stress}
% ----------------------------------------------------------------------------

The per-atom stress tensor $\boldsymbol{\sigma}_i$ is written in Voigt
notation: $({\sigma_{xx}}, {\sigma_{yy}}, {\sigma_{zz}},
{\sigma_{xy}}, {\sigma_{xz}}, {\sigma_{yz}})$.

The virial pressure is recovered from the per-atom contributions:
\begin{equation}
  P = -\frac{1}{3V} \sum_{i=1}^{N}
      \bigl(\sigma_{xx,i} + \sigma_{yy,i} + \sigma_{zz,i}\bigr)
  \label{eq:pressure}
\end{equation}

where $V = |\det\mathbf{H}|$ is the cell volume. This is written to the
comment line as \texttt{virial\_eV} for convenience.

% ----------------------------------------------------------------------------
\subsection{Image Counters}
\label{sec:xyzFull_images}

For PBC trajectories, each atom accumulates an image displacement counter
$(n_x, n_y, n_z) \in \mathbb{Z}^3$ tracking how many times it has crossed
each periodic boundary. The \emph{unwrapped} position is:

\begin{equation}
  \rvec_i^{\mathrm{unwrap}} = \rvec_i^{\mathrm{wrap}} + \mathbf{H}^\top \mathbf{n}_i
  \label{eq:unwrap}
\end{equation}

Unwrapped positions are required for correct MSD computation
(see \S\ref{sec:msd}).

% ----------------------------------------------------------------------------
\subsection{Comment Line (Full-State)}
\label{sec:xyzFull_comment}
% ----------------------------------------------------------------------------

The \texttt{.xyzFull} comment line carries all standard keys plus:

\begin{lstlisting}[caption={\texttt{.xyzFull} comment line.}]
step=1000 time_fs=50.000 energy_eV=-17641.320 temperature_K=300.2
  ke_eV=248.411 pe_eV=-17889.731 virial_eV=-52341.8
  pbc="T T T" lattice="22.624 0.0 0.0  0.0 22.624 0.0  0.0 0.0 22.624"
  hash=b7d3a941
\end{lstlisting}

Energy decomposition identity:
\begin{equation}
  E_{\mathrm{total}} = E_k + E_p,
  \qquad
  E_k = \sum_{i=1}^{N} \varepsilon_{k,i},
  \qquad
  E_p = \sum_{i=1}^{N} \varepsilon_{p,i}
  \label{eq:energy_decomp}
\end{equation}

% ============================================================================
\section{Bounding Box and Periodic Boundary Conditions}
\label{sec:pbc}
% ============================================================================

All format variants support the same PBC encoding in the comment line.
The encoding is activated when \texttt{pbc = "T T T"} (or any
\texttt{T}/\texttt{F} subset) appears alongside a \texttt{lattice} key.

% ----------------------------------------------------------------------------
\subsection{Comment-Line Encoding}
\label{sec:pbc_encoding}
% ----------------------------------------------------------------------------

\begin{lstlisting}[caption={Orthorhombic box.}]
pbc="T T T" lattice="22.624 0.0 0.0  0.0 22.624 0.0  0.0 0.0 22.624"
\end{lstlisting}

\begin{lstlisting}[caption={Triclinic box (monoclinic example).}]
pbc="T T T" lattice="11.312 0.0 0.0  5.656 9.797 0.0  0.0 0.0 22.624"
\end{lstlisting}

\begin{lstlisting}[caption={Slab geometry (periodic in x,y only).}]
pbc="T T F" lattice="22.624 0.0 0.0  0.0 22.624 0.0  0.0 0.0 60.000"
\end{lstlisting}

% ----------------------------------------------------------------------------
\subsection{Lattice Matrix}
\label{sec:pbc_vectors}
% ----------------------------------------------------------------------------

The nine values of the \texttt{lattice} key form the matrix $\mathbf{H}$
in row-major order:

\begin{equation}
  \mathbf{H} =
  \begin{bmatrix}
    \avec^\top \\
    \bvec^\top \\
    \cvec^\top
  \end{bmatrix}
  \in \mathbb{R}^{3 \times 3}
  \label{eq:lattice}
\end{equation}

An atom at fractional coordinates $\svec \in [0,1)^3$ has Cartesian position
$\rvec = \mathbf{H}^\top \svec$. The cell volume is $V = |\det \mathbf{H}|$.
For orthorhombic boxes, $\mathbf{H}$ is diagonal:
$L_x = a_{11}$, $L_y = a_{22}$, $L_z = a_{33}$.

% ----------------------------------------------------------------------------
\subsection{Reciprocal Lattice}
\label{sec:pbc_recip}
% ----------------------------------------------------------------------------

\begin{equation}
  \mathbf{H}^* = 2\pi \bigl(\mathbf{H}^\top\bigr)^{-1}
  \label{eq:recip}
\end{equation}

Rows of $\mathbf{H}^*$ are the reciprocal lattice vectors $\mathbf{a}^*,
\mathbf{b}^*, \mathbf{c}^*$ satisfying $\avec \cdot \mathbf{a}^* = 2\pi$,
$\avec \cdot \mathbf{b}^* = 0$, etc. Used internally by the Ewald summation
module; not written to XYZ files.

% ----------------------------------------------------------------------------
\subsection{Fractional Coordinate Conversion}
\label{sec:pbc_frac}
% ----------------------------------------------------------------------------

Cartesian to fractional:
\begin{equation}
  \svec = \bigl(\mathbf{H}^\top\bigr)^{-1} \rvec
  \label{eq:cart2frac}
\end{equation}

Wrapping a position into the primary cell:
\begin{equation}
  \svec_{\mathrm{wrap}} = \svec - \lfloor \svec \rfloor
  \label{eq:wrap}
\end{equation}

% ----------------------------------------------------------------------------
\subsection{Minimum Image Convention}
\label{sec:pbc_mic}
% ----------------------------------------------------------------------------

\begin{equation}
  \Delta\rvec_{ij}^{\mathrm{MIC}}
  = \Delta\rvec_{ij}
  - \mathbf{H}^\top
    \left\lfloor \mathbf{H}^{-\top} \Delta\rvec_{ij} + \tfrac{1}{2} \right\rfloor
  \label{eq:mic}
\end{equation}

For orthorhombic boxes:
\begin{equation}
  \Delta r_{\alpha,ij}^{\mathrm{MIC}}
  = \Delta r_{\alpha,ij}
    - L_\alpha \operatorname{round}\!\left(\frac{\Delta r_{\alpha,ij}}{L_\alpha}\right)
  \label{eq:mic_ortho}
\end{equation}

\texttt{pbc = "T T F"} applies PBC in $x$ and $y$ only (slab geometry).

% ============================================================================
\section{Bond Auto-Detection and Neighbor Lists}
\label{sec:bonds}
% ============================================================================

% ----------------------------------------------------------------------------
\subsection{Bond Detection Criterion}
\label{sec:bonds_criterion}
% ----------------------------------------------------------------------------

\begin{equation}
  \mathrm{bonded}(i,j)
  \iff
  d_{ij}^{\mathrm{MIC}}
  \;\leq\;
  \lambda \bigl(\rcov^{(i)} + \rcov^{(j)}\bigr)
  \label{eq:bond}
\end{equation}

Default $\lambda = 1.15$; overridable via
\texttt{[simulation].bond\_detect\_scale}.

% ----------------------------------------------------------------------------
\subsection{Neighbor List Construction}
\label{sec:bonds_neighborlist}
% ----------------------------------------------------------------------------

For $N > 500$ atoms, bond detection uses a cell-list neighbor structure.
Cut-off radius:
\begin{equation}
  \rcut = \lambda \max_{k} \bigl(2\,\rcov^{(k)}\bigr) + 0.3\,\ang
  \label{eq:rcut}
\end{equation}

Cells per axis:
\begin{equation}
  n_\alpha = \left\lfloor \frac{L_\alpha}{\rcut} \right\rfloor \geq 1
  \label{eq:ncells}
\end{equation}

Only atoms in the same cell or 26 neighboring cells need distance evaluation,
reducing complexity from $O(N^2)$ to $O(N)$.

% ----------------------------------------------------------------------------
\subsection{Coordination Number}
\label{sec:bonds_coordination}
% ----------------------------------------------------------------------------

\begin{equation}
  Z_i = \bigl|\{j \neq i : \mathrm{bonded}(i,j)\}\bigr|
  \label{eq:coordination}
\end{equation}

Mean coordination for species $s$:
\begin{equation}
  \bar{Z}_s = \frac{1}{N_s} \sum_{i : \mathrm{species}(i) = s} Z_i
  \label{eq:mean_coord}
\end{equation}

Expected values: NaCl $\bar{Z}=6$; diamond Si $\bar{Z}=4$;
graphene C $\bar{Z}=3$; BCC Fe $\bar{Z}=8$.

% ----------------------------------------------------------------------------
\subsection{Design Rationale}
\label{sec:bonds_rationale}
% ----------------------------------------------------------------------------

\begin{itemize}
  \item No pre-defined connectivity table is required.
  \item $\lambda$ is a single global knob. Ionic systems may need
    $\lambda \approx 1.3$; covalent systems use $\lambda = 1.15$.
  \item For metallic systems, bond detection is suppressed
    (\texttt{bond\_detect\_scale = 0.0}).
  \item For bead/polymer systems, topology is injected by the bead-builder.
  \item Covalent radius table in \texttt{include/vsim/element\_data.hpp}
    covers $Z = 1$--$118$.
\end{itemize}

% ============================================================================
\section{Trajectory Analysis}
\label{sec:analysis}
% ============================================================================

% ----------------------------------------------------------------------------
\subsection{Radial Distribution Function}
\label{sec:rdf}
% ----------------------------------------------------------------------------

\begin{equation}
  g(r) = \frac{V}{N^2}
         \left\langle
           \sum_{i} \sum_{j \neq i}
           \delta\!\left(r - d_{ij}^{\mathrm{MIC}}\right)
         \right\rangle
  \label{eq:rdf}
\end{equation}

Discretised into bins of width $\Delta r$:
\begin{equation}
  g(r_k) = \frac{V}{N^2} \cdot
            \frac{\langle n_k \rangle}{4\pi r_k^2 \Delta r}
  \label{eq:rdf_bin}
\end{equation}

Truncated at $r_{\max} = L_{\min}/2$ to avoid double-counting.
Species-resolved partial RDF:
\begin{equation}
  g_{AB}(r) = \frac{V}{N_A N_B}
              \left\langle
                \sum_{i \in A} \sum_{j \in B}
                \delta\!\left(r - d_{ij}^{\mathrm{MIC}}\right)
              \right\rangle
  \label{eq:rdf_partial}
\end{equation}

% ----------------------------------------------------------------------------
\subsection{Mean Squared Displacement}
\label{sec:msd}
% ----------------------------------------------------------------------------

Requires unwrapped positions $\rvec_i^{\mathrm{unwrap}}(t)$
from image counters in \texttt{.xyzFull}:

\begin{equation}
  \mathrm{MSD}(\tau) =
  \frac{1}{N}
  \left\langle
    \sum_{i=1}^{N}
    \left\|
      \rvec_i^{\mathrm{unwrap}}(t_0 + \tau) - \rvec_i^{\mathrm{unwrap}}(t_0)
    \right\|^2
  \right\rangle_{t_0}
  \label{eq:msd}
\end{equation}

Solid: $\mathrm{MSD}(\tau) \to C$ (bounded).
Diffusing liquid: $\mathrm{MSD}(\tau) \sim 6D\tau$ where $D$ is the
self-diffusion coefficient.

% ----------------------------------------------------------------------------
\subsection{Velocity Autocorrelation and VDOS}
\label{sec:vacf}
% ----------------------------------------------------------------------------

Normalised VACF:
\begin{equation}
  C_v(\tau) =
  \frac{
    \left\langle
      \sum_{i} \mathbf{v}_i(t_0+\tau) \cdot \mathbf{v}_i(t_0)
    \right\rangle_{t_0}
  }{
    \left\langle
      \sum_{i} \|\mathbf{v}_i(t_0)\|^2
    \right\rangle_{t_0}
  }
  \label{eq:vacf}
\end{equation}

Vibrational density of states (VDOS):
\begin{equation}
  g(\nu) = \int_{-\infty}^{\infty} C_v(\tau)\, e^{-2\pi i \nu \tau}\, d\tau
  \label{eq:vdos}
\end{equation}

Requires velocities; \texttt{.xyz} and \texttt{.xyzf} cannot be used.

% ============================================================================
\section{Provenance Hashing}
\label{sec:provenance}
% ============================================================================

\subsection{Frame Hash}
\label{sec:provenance_frame}

\begin{equation}
  H_{\mathrm{frame}} = \mathrm{SHA256}\bigl(
    \texttt{step} \;\|\; \texttt{n\_atoms}
    \;\|\; \texttt{species} \;\|\; \texttt{positions}
  \bigr)[0:8]
  \label{eq:frame_hash}
\end{equation}

Position values are rounded to 6 decimal places before hashing.

\subsection{File Hash}
\label{sec:provenance_file}

\begin{equation}
  H_{\mathrm{file}} = \mathrm{SHA256}\bigl(\mathrm{file}[0:10240]\bigr)
  \label{eq:file_hash}
\end{equation}

Recomputed on restart; mismatch is a hard error.

\subsection{Hash Chain}
\label{sec:provenance_chain}

\begin{equation}
  H_k = \mathrm{SHA256}\bigl(H_{k-1} \,\|\, H_{\mathrm{frame},k}\bigr)[0:8]
  \label{eq:hash_chain}
\end{equation}

Opt-in via \texttt{[kernel].symbolic\_trace = true}; stored in
\texttt{write\_pipeline\_audit\_jsonl}.

% ============================================================================
\section{Per-Format Validation}
\label{sec:validation}
% ============================================================================

\subsection{Universal Checks}
\begin{enumerate}
  \item Atom count on line~1 must equal the number of atom lines per frame.
  \item All coordinate values must be finite (no \texttt{NaN}/\texttt{Inf}).
  \item Species column must not be empty.
  \item \texttt{step} $\geq 0$; \texttt{time\_fs} $\geq 0$.
\end{enumerate}

\subsection{\texttt{.xyza} Checks}
\begin{enumerate}
  \item Charge model must not change between frames.
  \item Velocity and force values must be finite.
  \item $\|\mathbf{v}_i\| < 100\,\ang\,\fs^{-1}$ (warning above this).
  \item $\|\mathbf{F}_i\| < 10^4\,\eV\,\ang^{-1}$ (warning above this).
\end{enumerate}

\subsection{\texttt{.xyzc} Checks}
\begin{enumerate}
  \item Matching \texttt{BEGIN}/\texttt{END} sentinels required.
  \item \texttt{n\_atoms} must equal the first frame's atom count.
  \item \texttt{file\_hash} must match recomputed value on restart.
  \item \texttt{format\_version} must be $\leq$ reader's maximum.
\end{enumerate}

\subsection{\texttt{.xyzFull} Checks}
\begin{enumerate}
  \item $|E_{\mathrm{total}} - (E_k + E_p)| < 10^{-6}\,\eV$.
  \item $\varepsilon_{k,i} \geq 0$ for all atoms.
  \item Stress Voigt components self-consistent.
  \item $|\Delta n_\alpha| \leq 1$ per atom per axis between consecutive frames.
\end{enumerate}

% ============================================================================
\section{Interoperability}
\label{sec:interop}
% ============================================================================

\subsection{ASE}
ASE \texttt{extxyz} is compatible with \texttt{.xyza} when the
\texttt{Properties} key is present. Column mapping:
\texttt{pos} $\leftrightarrow$ $x,y,z$;
\texttt{charge} $\leftrightarrow$ $q$;
\texttt{vel} $\leftrightarrow$ $v_x,v_y,v_z$;
\texttt{forces} $\leftrightarrow$ $F_x,F_y,F_z$.
Stress and image counters are not mapped.

\subsection{VESTA}
VESTA reads standard \texttt{.xyz}. Export with
\texttt{XyzFormat::Xyz} and \texttt{suppress\_props = true}.

\subsection{Ovito}
Ovito reads \texttt{.xyz} and \texttt{.xyza} via its extended-XYZ parser.
It accepts both uppercase \texttt{Lattice} and lowercase \texttt{lattice};
VSEPR-SIM writes lowercase.

% ============================================================================
\section{Export Profile Connection}
\label{sec:export}
% ============================================================================

\begin{center}
\begin{tabular}{lll}
  \toprule
  \texttt{[export]} flag & File written & Format \\
  \midrule
  \texttt{write\_xyz = true}        & \texttt{<prefix>.xyz}            & \texttt{.xyz} \\
  \texttt{write\_xyzf = true}       & \texttt{<prefix>.xyzf}           & \texttt{.xyzf} \\
  \texttt{write\_xyzFull = true}    & \texttt{<prefix>.xyzFull}        & \texttt{.xyzFull} \\
  \texttt{write\_checkpoint = true} & \texttt{<prefix>\_step\{N\}.xyzc}& \texttt{.xyzc} \\
  \bottomrule
\end{tabular}
\end{center}

Export profile expansion:
\texttt{minimal} $\to$ \texttt{.xyz} only;
\texttt{standard} $\to$ \texttt{.xyz}, \texttt{.xyzf};
\texttt{research\_report} $\to$ adds \texttt{.xyzFull};
\texttt{publication} $\to$ all + \texttt{.xyzc} every 1000 steps.

% ============================================================================
\section{Implementation Mapping}
\label{sec:impl}
% ============================================================================

\begin{table}[h]
\centering
\begin{tabular}{lll}
  \toprule
  Source file & Responsibility & Format(s) \\
  \midrule
  \texttt{xyz\_io.hpp / .cpp}         & Core reader/writer           & \texttt{.xyz}, \texttt{.xyza} \\
  \texttt{xyz\_full.hpp / .cpp}       & Full-state reader/writer     & \texttt{.xyzFull} \\
  \texttt{xyz\_checkpoint.hpp / .cpp} & Checkpoint I/O               & \texttt{.xyzc} \\
  \texttt{xyz\_formation.hpp / .cpp}  & Formation trajectory I/O     & \texttt{.xyzf} \\
  \texttt{element\_data.hpp}          & Covalent radius \& mass table & all \\
  \texttt{pbc\_utils.hpp}             & MIC, wrap/unwrap, cell volume & all \\
  \texttt{cell\_list.hpp}             & Neighbor-list construction    & all \\
  \texttt{bond\_detect.hpp}           & Bond auto-detection           & all \\
  \texttt{traj\_analysis.hpp / .cpp}  & RDF, MSD, VACF, VDOS          & \texttt{.xyza}, \texttt{.xyzFull} \\
  \texttt{provenance.hpp}             & Frame and file hash utilities  & all \\
  \texttt{vsim\_export.hpp / .cpp}    & Export dispatch layer         & all \\
  \bottomrule
\end{tabular}
\caption{Source files implementing the XYZ format family.}
\label{tab:impl}
\end{table}

\subsection{Core Frame Struct}

\begin{lstlisting}[language=C++, caption={Core frame struct (xyz\_io.hpp).}]
namespace vsim {

struct XyzFrame {
    int                       n_atoms;
    XyzFormat                 format;
    CommentFields             comment;
    std::vector<std::string>  species;
    std::vector<Vec3>         positions;     // Ang
    std::vector<double>       charges;       // e
    std::vector<Vec3>         velocities;    // Ang/fs   (.xyza+)
    std::vector<Vec3>         forces;        // eV/Ang   (.xyza+)
    // .xyzFull only:
    std::vector<double>       ke_per_atom;   // eV
    std::vector<double>       pe_per_atom;   // eV
    std::vector<Stress3>      stress;        // eV/Ang^3, Voigt
    std::vector<IVec3>        image_counts;  // PBC image displacement
};

} // namespace vsim
\end{lstlisting}

\subsection{Reader API}

\begin{lstlisting}[language=C++, caption={XYZ reader API.}]
namespace vsim {

bool read_frame(std::istream& in, XyzFrame& out,
                XyzFormat fmt = XyzFormat::Auto);

std::vector<XyzFrame> read_all(
    const std::filesystem::path& p,
    XyzFormat fmt = XyzFormat::Auto);

/// Seek to frame k in a .xyzc file (O(1) line offset).
bool seek_frame(std::istream& in, int k,
                int n_atoms, int n_header_lines);

} // namespace vsim
\end{lstlisting}

\subsection{Writer API}

\begin{lstlisting}[language=C++, caption={XYZ writer API.}]
namespace vsim {

struct WriteOptions {
    int  precision_pos   = 6;
    int  precision_other = 4;
    bool suppress_props  = false;
    bool write_hash      = true;
    bool chain_hash      = false;
};

void write_frame(std::ostream& out, const XyzFrame& frame,
                 XyzFormat fmt, const WriteOptions& opts = {});

void append_frame(const std::filesystem::path& p,
                  const XyzFrame& frame,
                  XyzFormat fmt, const WriteOptions& opts = {});

} // namespace vsim
\end{lstlisting}

\subsection{Analysis API}

\begin{lstlisting}[language=C++, caption={Trajectory analysis API (traj\_analysis.hpp).}]
namespace vsim {

struct RdfResult {
    std::vector<double> r;
    std::vector<double> g_r;
    double              dr;
};

struct MsdResult {
    std::vector<double> tau;
    std::vector<double> msd;
    std::string         regime; // solid_bounded | diffusive | mixed
    double              D;      // diffusion coeff (Ang^2/fs)
};

RdfResult compute_rdf(const std::vector<XyzFrame>& traj,
                      double dr = 0.05, double r_max = 0.0);

MsdResult compute_msd(const std::vector<XyzFrame>& traj);

} // namespace vsim
\end{lstlisting}

% ============================================================================
\section{Typical Energy and Observable Ranges}
\label{sec:ranges}
% ============================================================================

\begin{center}
\begin{tabular}{llllll}
  \toprule
  System & $N$ & $E_{\mathrm{total}}$ (eV) & $T$ (K) & $d_{\mathrm{NN}}$ (\AA) & $\bar{Z}$ \\
  \midrule
  H$_2$O (1 molecule)              & 3   & $-14$ to $-13$         & 0--600  & 0.96       & O:2 \\
  NaCl rocksalt ($4{\times}4{\times}4$)& 512 & $-17700$ to $-17400$ & 0--1200 & 2.70--2.85 & 6 \\
  Si diamond ($3{\times}3{\times}3$)   & 216 & $-900$ to $-860$     & 0--1600 & 2.33--2.36 & 4 \\
  Graphene (96 C)                  & 96  & $-550$ to $-530$       & 0--2000 & 1.40--1.43 & 3 \\
  Fe BCC ($4{\times}4{\times}4$)   & 128 & $-820$ to $-800$       & 0--1800 & 2.45--2.50 & 8 \\
  UO$_2$ fluorite ($2{\times}2{\times}2$)& 96 & $-1900$ to $-1800$ & 0--3000 & 2.34       & U:8 \\
  Al FCC ($4{\times}4{\times}4$)   & 256 & $-860$ to $-830$       & 0--900  & 2.85--2.87 & 12 \\
  TiO$_2$ rutile ($2{\times}2{\times}4$)& 192& $-2800$ to $-2700$  & 0--2000 & 1.94       & Ti:6 \\
  \bottomrule
\end{tabular}
\end{center}

\medskip
\noindent\textbf{Energy.} Per simulation cell, not per atom. Per-atom values
in \texttt{.xyzFull} and \texttt{write\_analysis\_json}.

\medskip
\noindent\textbf{Temperature.} Instantaneous kinetic temperature:
\begin{equation}
  T = \frac{2 E_k}{3 N \kB}
  \label{eq:T_kin}
\end{equation}

Fluctuates substantially for small $N$. Block-averaged temperature in
\texttt{write\_metrics\_tsv}.

\medskip
\noindent\textbf{MSD plateaus (solid, 300\,K):}
NaCl $0.05$--$0.15\,\ang^2$;
Si $0.01$--$0.05\,\ang^2$;
Fe $0.02$--$0.08\,\ang^2$.
Values substantially above these suggest melting or a runaway trajectory.

\vfill
\begin{center}
  \rule{0.4\textwidth}{0.4pt}\\[0.5em]
  \textit{XYZ File Format Family --- VSEPR-SIM Theoretical Division.
  Plain text was good enough for Kohn and Sham; it will do.}
\end{center}

\end{document}
